\section{Background and Replicated Claims}
\label{sec:background}

\subsection{Sampling configuration spaces}

A configurable system with $k$ options spans a configuration space whose
valid members $\mathcal{V}$ are constrained by a variability model
\citep{apel2013feature}. A \emph{sampling strategy} selects a subset
$S \subseteq \mathcal{V}$, $|S| = n$, whose measurements serve a
downstream purpose: training a performance model
\citep{guo2013variability,siegmund2015performance,ha2019deepperf},
finding a near-optimal configuration \citep{oh2017srs,nair2020flash}, or,
in our setting, bounding worst-case (or best-case) performance. Strategy
families include uniform random sampling (in practice often
solver-mediated \citep{plazar2019uniform,heradio2022uniform}), coverage
oriented $t$-wise sampling \citep{medeiros2016comparison,baranov2024baital,hess2024unwise},
distance- and diversity-based sampling \citep{kaltenecker2019distance},
and adaptive strategies that use intermediate measurements to steer
selection \citep{sarkar2015cost,nair2017bad}.

\emph{Statistical Recursive Search} (\srs) \citep{oh2017srs}, the
strategy recommended by \citet{weis2026risk}, is adaptive: it repeatedly
draws a small random batch (25 configurations), identifies option
decisions shared by the two best configurations of the batch, keeps the
decisions whose enabled/disabled performance difference is significant
under Welch's $t$-test ($p<0.05$), restricts the space accordingly, and
recurses; if recursion stops before the budget is exhausted, the
remainder is filled with random configurations.

\subsection{Posterior risk and posterior confidence}
\label{sec:posterior}

Let $y(c)$ be the measured target of configuration $c \in \mathcal{V}$
and let $S$ be a sample. For a worst-case (maximum) bound, the sampled
bound is $\hat{b} = \max_{c \in S} y(c)$, and the \emph{posterior risk}
is the fraction of the reference universe that violates it:
\begin{equation}
  R(S) \;=\; \frac{\big|\{\,c \in \mathcal{V} : y(c) > \hat{b}\,\}\big|}{|\mathcal{V}|}.
  \label{eq:risk}
\end{equation}
Repeating the sampling $r$ times yields an empirical distribution of
risks; the \emph{posterior confidence} at tolerance $\varepsilon$ is the
fraction of repetitions with $R \le \varepsilon$, i.e., the empirical CDF
of risk. Integrating this CDF over $\varepsilon \in [0,1]$ gives the
\emph{integrated posterior confidence} $A_{\mathrm{post}} = 1 -
\overline{R}$, a scalar summary (higher is better). Scenario theory
provides a strategy-agnostic \emph{prior} baseline: for uniform random
samples of size $N$, confidence at tolerance $\varepsilon$ is
$1-(1-\varepsilon)^{N+1}$ (we follow the artifact's exponent convention
and report a sensitivity with exponent $N$). The \emph{integrated
confidence discrepancy} $\omega = A_{\mathrm{post}} - A_{\mathrm{prior}}$
measures how much better (or worse) a strategy's guarantees are than the
random-sampling prior at the same size. All quantities are computed from
measurements only; no learner is involved.

\citet{weis2026risk} evaluate five strategies (uniform random,
solver-based, diversified distance-based, $t$-wise, \srs) on seven
systems with exhaustively measured configuration spaces (7z, BerkeleyDB-C,
Dune, Hipacc, JavaGC, LLVM, Polly), at a per-system budget $t2$ matching
the size of a pairwise-coverage sample (56--843 configurations), with 100
repetitions. We refer to their headline results as four claims:

\begin{description}
\item[C1 (discriminative power).] Posterior risk and confidence
  distinguish sampling strategies at the same budget, enabling rigorous
  comparison beyond observed performance values.
\item[C2 (\srs{} superiority).] \srs{} yields consistently tighter
  guarantees---lower posterior risk, higher posterior confidence---than
  the alternatives at the same budget.
\item[C3 (solver unreliability).] Plain solver-based sampling is the
  least reliable strategy for probabilistic guarantees.
\item[C4 (budget scaling).] Guarantees improve as the sample size grows.
\end{description}

\subsection{Sample size versus sampling cost}
\label{sec:costargument}

The budget in C1--C4 is a \emph{count}. The construct-validity concern is
that counts are not effort: measuring configuration $c$ costs some
$\kappa(c) > 0$ (seconds of wall-clock, joules, dollars), and strategies
select different configurations. We quantify the gap with the \emph{cost
share}
$\%cost(S) = \sum_{c\in S}\kappa(c) / \sum_{c\in\mathcal{V}}\kappa(c)$,
the \emph{size share} $\%n(S) = |S|/|\mathcal{V}|$, and their ratio, the
\emph{calibration ratio} $\%cost/\%n$ (1.0 iff the sample is
cost-representative). Cost-indexing the prior baseline is then immediate:
a strategy that spent cost share $\%cost$ has an \emph{effective
cost-equivalent size} $N_{\mathrm{eff}} = \%cost \cdot |\mathcal{V}|$,
and $\omega_{\mathrm{cost}} = A_{\mathrm{post}} -
A_{\mathrm{prior}}(N_{\mathrm{eff}})$ asks whether its guarantees beat
what uniform random sampling would be expected to deliver \emph{for the
same spend}. Analogously, Pareto frontiers over (confidence, budget) can
use either $n$ or measured cost as the budget axis; comparing the two
frontiers (by membership and Jaccard overlap) quantifies how much the
choice of cost construct changes the decision surface. Throughout,
$\kappa$ is itself a measured property of the corpus (runtime or energy),
not a model.

\subsection{Replication terminology}

Following \citet{shull2008role} and \citet{gomez2014understanding}, \rqone{}
is a \emph{close} replication on the original artifact (same data, same
strategy samples, independently reimplemented analysis) and a
\emph{conceptual} replication on five further corpora (same constructs
and hypotheses, different populations, partially different strategy
sets). \rqtwo{} is an \emph{extension}: it changes one construct---the
budget axis---while holding everything else fixed.
